\section{Mathematical Framework}
\label{sec:mathematical_framework}

This section presents the mathematical framework for the stochastic harmonic hosting capacity problem of a phase-balanced three-phase power system, considering uncertainty on the angle of the harmonic current injected by the connected harmonic units. To this end, the formulation for the deterministic harmonic hosting capacity presented in~\cite{vanacker_hosting_2026} is extended. For brevity, the original formulation is not repeated in this section, rather the relevant variables and constraints are highlighted and transformed to their stochastic counterpart. %

\subsection{Background on Polynomial Chaos Expansion}
\label{sec:polynomial_chaos_expansion}

This section introduces the necessary background on polynomial chaos expansion, a method for representing random variables in terms of polynomial functions. Concretely, a vector of input random variables, e.g., harmonic current angles, is propagated through the set of constraints of the harmonic hosting capacity problem to obtain a vector of unknown output random variables, e.g., harmonic bus voltages. For a detailed treatment of the subject, the reader is referred to~\cite{sullivan_introduction_2015}.

\subsubsection{Uncertainty Quantification}
%It is assumed that the stochastic input is captured by a real-valued finite-variance stochastic germ~$\bar{\omega} = [\omega_{1},...,\omega_{N}]^{\text{T}}$ of independent random variables. Optimization solvers are unable to handle random variables directly, necessitating a deterministic equivalent of the stochastic problem. For this purpose, polynomial chaos expansion approximates any real-valued random variable~$\RVX$ of finite variance that is a function of the germ~$\bar{\omega}$ as a linear combination~$\ApproxRVX$ of the orthogonal polynomial basis~$[\Polynoom_{\PCEIndex}(\bar{\omega})]_{\PCEIndex \in \PCESet}$, where $\PCESet\subseteq\NaturalNumberSet_{0}$, as,
It is assumed that the stochastic input is captured by a real-valued finite-variance stochastic germ~$\bar{\omega} = [\omega_{1},...,\omega_{N}]^{\text{T}}$ of independent random variables\changed{, i.e., the inputs from which all random quantities of the problem are derived}. Optimization solvers are unable to handle random variables directly, necessitating a deterministic equivalent of the stochastic problem. For this purpose, polynomial chaos expansion approximates any real-valued random variable~$\RVX$ of finite variance that is a function of the germ~$\bar{\omega}$ as a linear combination~$\ApproxRVX$ of the orthogonal polynomial basis~$[\Polynoom_{\PCEIndex}(\bar{\omega})]_{\PCEIndex \in \PCESet}$, where $\PCESet\subseteq\NaturalNumberSet_{0}$, as,
\begin{IEEEeqnarray}{ r C l }
    \ApproxRVX &=& \sum_{\PCEIndex \in \PCESet} \CoefficientX_{\PCEIndex} \cdot \Polynoom_{\PCEIndex} =
    \sum_{\PCEIndex \in \PCESet} \frac{\langle \RVX, \Polynoom_{\PCEIndex} \rangle}{\langle \Polynoom_{\PCEIndex}, \Polynoom_{\PCEIndex} \rangle} \cdot \Polynoom_{\PCEIndex}, \label{eq:gpc_expansion}
\end{IEEEeqnarray}
where~$\CoefficientX_{\PCEIndex} \in \RealNumberSet$ denote the polynomial chaos expansion coefficients and $\langle \cdot,\cdot \rangle$ denotes the inner product with respect to the probability measure of the germ. Following the Askey scheme~\cite{sullivan_introduction_2015}, each Beta-distributed random variable corresponds to the Jacobi family of orthogonal polynomials; the multivariate basis is the tensor product of the univariate ones, truncated to total degree~$\delta$, and is constructed numerically with \textsc{PolyChaos.jl}~\cite{muhlpfordt_polychaos_2020}.

\subsubsection{Uncertainty Projection onto Equality Constraints}
Through polynomial chaos expansion, any stochastic equality constraint is equivalently represented by $|\PCESet|$ deterministic equality constraints. The cardinality of the model~$|\PCESet|$ is determined by the dimension of the germ and the highest polynomial degree as,
\begin{IEEEeqnarray}{ r C l }
    |\PCESet| &=& \frac{(|\bar{\omega}| + \delta)!}{|\bar{\omega}|!\,\delta!}, \label{eq:cardinality}
\end{IEEEeqnarray}
where~$\delta$ denotes the maximum degree of the polynomial basis functions.

The Galerkin projection forms a system of equations for the stochastic modes~(Chapter~12 in~\cite{sullivan_introduction_2015}). For the harmonic hosting capacity problem, two operations stand out: summation and multiplication. For three approximated random variables~$\ApproxRVX$,~$\ApproxRVY$ and~$\ApproxRVZ$ defined on the same basis, the Galerkin projection, denoted by~$\gp$, of both operations is,
\begin{IEEEeqnarray}{ l C l l}
    \ApproxRVZ = \ApproxRVX + \ApproxRVY &\gp& \CoefficientZ_{\PCEIndex} = \CoefficientX_{\PCEIndex} + \CoefficientY_{\PCEIndex},~&\forall \PCEIndex \in \PCESet,~ \label{eq:gp_sum} \\
    \ApproxRVZ = \ApproxRVX \cdot \ApproxRVY &\gp& \CoefficientZ_{\PCEIndex} = \mkern-15mu \sum_{\bar{\kappa} \in P^{*}(\PCESet,2)} \mkern-15mu M_{\bar{\kappa},\PCEIndex} \,(\CoefficientX_{\bar{\kappa}_{1}} \cdot \CoefficientY_{\bar{\kappa}_{2}}),\quad&\forall \PCEIndex \in \PCESet,~ \label{eq:gp_product}
\end{IEEEeqnarray}
where~$P^{*}(\PCESet,2)$ denotes the 2-tuples of~$\PCESet$ and the multiplication tensor,
\begin{IEEEeqnarray}{ r C l }
    M_{\bar{\kappa},\PCEIndex} = \frac{\langle \Polynoom_{\bar{\kappa}_{1}}\Polynoom_{\bar{\kappa}_{2}},\Polynoom_{k} \rangle}{\langle \Polynoom_{k},\Polynoom_{k}\rangle}, \label{eq_m_tensor_def}
\end{IEEEeqnarray}
is computed ahead of time. For $\PCEIndex=0$, orthogonality implies that $M_{\bar{\kappa},0}$ is diagonal with strictly positive entries; for $\PCEIndex \neq 0$, the tensor has mixed-sign entries, and the matrix $[M_{\bar{\kappa},k}]_{\bar\kappa}$ is generally indefinite.

\subsubsection{Chance-Constraint Reformulation of Inequality Constraints}
All harmonic limits in this work are upper bounds. A chance constraint ensures that the probability of~$\RVX$ exceeding a deterministic limit~$x^{\text{max}}$ is below a risk level~$\varepsilon$, i.e., $\Probability(\RVX \leq x^{\text{max}}) \geq 1 - \varepsilon$. It is enforced through the moment-based reformulation,
\begin{IEEEeqnarray}{ r C l }
    \Expectation(\RVX) + \lambda(\varepsilon)\sqrt{\Variance(\RVX)} \leq x^{\text{max}}, \label{eq:cc_constraint}
\end{IEEEeqnarray}
with the expectation and variance of an approximated random variable~$\ApproxRVX$ given in terms of its coefficients as,
\begin{IEEEeqnarray}{ r C l }
    \Expectation[\ApproxRVX]   &=& \CoefficientX_{0}, \label{eq:gpc_expectation} \\
    \Variance[\ApproxRVX]      &=& \mkern-10mu \sum_{\PCEIndex \in \PCESet \backslash \{0\}} \mkern-10mu \langle \Polynoom_{\PCEIndex}, \Polynoom_{\PCEIndex} \rangle \, \CoefficientX_{\PCEIndex}^{2}, \label{eq:gpc_variance}
\end{IEEEeqnarray}
so that~\eqref{eq:cc_constraint} is a second-order cone constraint in the coefficients~$\{\CoefficientX_{\PCEIndex}\}_{\PCEIndex \in \PCESet}$. The risk factor~$\lambda(\varepsilon)$ is calibrated per constraint as $\lambda^{*} = (q_{1-\varepsilon}(\RVX)-\Expectation[\ApproxRVX])/\sqrt{\Variance[\ApproxRVX]}$, with the $(1-\varepsilon)$ quantile~$q_{1-\varepsilon}$ of samples of the constrained quantity and the moments~\eqref{eq:gpc_expectation}--\eqref{eq:gpc_variance} of its expansion, and is updated with the hosting capacity until convergence. As it is calibrated on the moments of the expansion, the bound equals the sampled quantile irrespective of the truncation of the basis, i.e., it targets the risk level~$\varepsilon$ exactly.

\subsection{Uncertainty on Angle of Harmonic Current Injection}
\label{sec:harmonic_current_angle_uncertainty}

%This section discusses the treatment of the uncertain angles~$\phi_{\Unit,\Harmonic}$ of the harmonic currents injected by the harmonic units~$\Unit \in \HarmonicUnitSet$, for all harmonics~$\Harmonic \in \HarmonicSet$. Throughout, a polynomial basis of degree~$\delta = 2$ is used, which captures the non-linear nature of the power flow equations without significant loss of accuracy~\cite{muhlpfordt_chance_2019}. For narrow angle ranges, it keeps at least~92\,\% of the variance of the cosine and sine of the angle up to the 13th harmonic, but only~39\,\% at the 25th harmonic of the thyristor bridges. The calibration of the risk factor on the moments of the expansion, Section~\ref{sec:polynomial_chaos_expansion}, compensates for this truncation.
This section discusses the treatment of the uncertain angles~$\phi_{\Unit,\Harmonic}$ of the harmonic currents injected by the harmonic units~$\Unit \in \HarmonicUnitSet$, for all harmonics~$\Harmonic \in \HarmonicSet$. Throughout, a polynomial basis of degree~$\delta = 2$ is used, which captures the non-linear nature of the power flow equations without significant loss of accuracy~\cite{muhlpfordt_chance_2019}. \changed{For the angle distributions of Section~\ref{sec:uncertainty}, computed by quadrature,} it keeps at least~92\,\% of the variance of the cosine and sine of the angle up to the 13th harmonic, but only~39\,\% at the 25th harmonic of the thyristor bridges. The calibration of the risk factor on the moments of the expansion, Section~\ref{sec:polynomial_chaos_expansion}, compensates for this truncation.

\subsubsection{Random Variable per Harmonic Unit}
As shown in Section~\ref{sec:uncertainty}, the phase angles of a converter across the harmonic orders are strongly dependent, and their marginal distributions are well approximated by a single shared Beta distribution with harmonic-specific location and scale. Each harmonic unit~$\Unit$ is therefore assigned one random variable~$\omega_{\Unit} \sim \mathcal{B}(\hat{\alpha},\hat{\beta})$ on~$[0,1]$, which drives its angle at every harmonic through an affine map,
\begin{IEEEeqnarray}{ r C l }
    \phi_{\Unit,\Harmonic} &=& v^{\text{min}}_{\Harmonic} + (v^{\text{max}}_{\Harmonic} - v^{\text{min}}_{\Harmonic})\,\omega_{\Unit} + \Harmonic\,\theta_{\Node(\Unit)}, \label{eq:angle_affine}
\end{IEEEeqnarray}
while the random variables of different units are independent. At the higher harmonics, the angle range may exceed~$2\pi$, e.g., 660$^{\circ}$ at the 25th harmonic of the thyristor bridges. The affine map then wraps around the circle, which leaves the phasor~$e^{\im\phi_{\Unit,\Harmonic}}$ and the dependence between the harmonics of a unit well defined. For a mix of converter technologies, the shape parameters~$\hat{\alpha}$, $\hat{\beta}$ and the angle ranges~$[v^{\text{min}}_{\Harmonic}, v^{\text{max}}_{\Harmonic}]$ are those of the technology of unit~$\Unit$, Section~\ref{sec:uncertainty}, and the basis is the tensor product of the univariate bases of the units.

\subsubsection{Local Angle Reference}
The angle ranges of Section~\ref{sec:uncertainty} result from a converter model driven by its own terminal voltage, i.e., they are defined relative to the fundamental voltage at the bus~$\Node(\Unit)$ the unit connects to. A shift of the fundamental by its angle~$\theta_{\Node}$ rotates the $\Harmonic$-th harmonic phasor by~$\Harmonic\,\theta_{\Node}$, which gives the last term of~\eqref{eq:angle_affine}. The fundamental angles follow from the fundamental power flow, which is solved first. Since~$\theta_{\Node}$ depends on the operating point and not only on the network topology, units of the same type at different buses inject at different absolute angles: the local angle reference is a physically required source of angle diversity, and it is the reference used throughout this work.

\subsubsection{Exact Projection of the Trigonometric Functions}
In the deterministic formulation~\cite{vanacker_hosting_2026}, the harmonic current magnitude~$\CurrentMagnitude_{\Unit,\Harmonic}$ is related to its rectangular components through the angle~$\phi_{\Unit,\Harmonic}$,
\begin{IEEEeqnarray}{ l }
        \cos(\phi_{\Unit,\Harmonic}) \, \CurrentMagnitude_{\Unit,\Harmonic} =
        \RealCurrent_{\Unit,\Harmonic}, \quad \forall \Unit \in \HarmonicUnitSet, \Harmonic \in \HarmonicSet,\label{eq:real_harmonic_unit_current} \\
        \sin(\phi_{\Unit,\Harmonic}) \, \CurrentMagnitude_{\Unit,\Harmonic} =
        \ImaginaryCurrent_{\Unit,\Harmonic}, \quad \forall \Unit \in \HarmonicUnitSet, \Harmonic \in \HarmonicSet. \label{eq:imag_harmonic_unit_current}
\end{IEEEeqnarray}
Since the angle~\eqref{eq:angle_affine} is affine in a single random variable, $\cos(\phi_{\Unit,\Harmonic})$ and $\sin(\phi_{\Unit,\Harmonic})$ are functions of~$\omega_{\Unit}$ alone, and their coefficients follow from~\eqref{eq:gpc_expansion}, for every index~$\PCEIndex$ supported on~$\omega_{\Unit}$ alone and zero otherwise,
\begin{IEEEeqnarray}{ r C l }
    \cos(\phi_{\Unit,\Harmonic})_{\PCEIndex} &=& \frac{\langle \cos(\phi_{\Unit,\Harmonic}), \Polynoom_{\PCEIndex} \rangle}{\langle \Polynoom_{\PCEIndex},\Polynoom_{\PCEIndex}\rangle}, \label{eq:trig_quadrature}
\end{IEEEeqnarray}
and analogously for the sine. %
%The coefficients~$\cos(\phi_{\Unit,\Harmonic})_{\PCEIndex}$ and~$\sin(\phi_{\Unit,\Harmonic})_{\PCEIndex}$ are parameters, not optimization variables.
The coefficients~$\cos(\phi_{\Unit,\Harmonic})_{\PCEIndex}$ and~$\sin(\phi_{\Unit,\Harmonic})_{\PCEIndex}$ are parameters, not optimization variables\changed{, and therefore do not affect the convexity of the optimization problem}.

\subsection{General Stochastic Harmonic Hosting Capacity}
\label{sec:stochastic_harmonic_hosting_capacity}

The stochastic model determines a deterministic harmonic hosting capacity given uncertainty on the harmonic current angles, while complying with chance constraints on (harmonic) voltages and currents. Given that the problem is formulated from the perspective of the power system operator, deterministic harmonic budgets are the practical outcome for its customers. A deterministic hosting capacity implies that~$\CurrentMagnitude_{\Unit,\Harmonic,\PCEIndex} = 0$ for any non-zero index~$\PCEIndex$, so the counterparts of~\eqref{eq:real_harmonic_unit_current} and~\eqref{eq:imag_harmonic_unit_current} are,
\begin{IEEEeqnarray}{ l }
    \cos(\phi_{\Unit,\Harmonic})_{\PCEIndex} \, \CurrentMagnitude_{\Unit,\Harmonic,0} = \RealCurrent_{\Unit,\Harmonic,\PCEIndex},
        \quad \forall \Unit \in \HarmonicUnitSet, \Harmonic \in \HarmonicSet, \PCEIndex \in \PCESet,
    \label{eq:real_harmonic_unit_current_pce} \\
    \sin(\phi_{\Unit,\Harmonic})_{\PCEIndex} \, \CurrentMagnitude_{\Unit,\Harmonic,0} = \ImaginaryCurrent_{\Unit,\Harmonic,\PCEIndex},
        \quad \forall \Unit \in \HarmonicUnitSet, \Harmonic \in \HarmonicSet, \PCEIndex \in \PCESet,
    \label{eq:imag_harmonic_unit_current_pce}
\end{IEEEeqnarray}
which are linear in the decision variables. Any variables related to the fairness objective are also deterministic. All other parameters, e.g., network impedances, are deterministic.

\subsubsection{Reformulation of Linear Constraints}
\label{sec:reformulation_of_linear_constraints}
The network equations of the current-voltage formulation in~\cite{vanacker_hosting_2026}, i.e., the reference bus, Kirchhoff's current law, Ohm's law, branch current and transformer constraints~(10)--(31), are linear, as are the fairness constraints~(34), (36) and~(38)--(39) and the harmonic current injection constraints~(40)--(41), of which the latter are treated above. Supported by~\eqref{eq:gp_sum}, a generic linear constraint $a\,\RVX + b\,\RVY = c$ with deterministic parameters is repeated over all indices, $a\,\RVX_{k} + b\,\RVY_{k} = c_{k}$, $\forall \PCEIndex \in \PCESet$, where $c_{k} = c$ for $k=0$ and $c_{k}=0$ otherwise. Since each voltage and current coefficient is then linear in the injected current coefficients, the physical variables are non-zero only on the indices supported by a single random variable%
, which reduces the number of variables.

\subsubsection{Reformulation of Quadratic Constraints and Problem Class}
\label{sec:reformulation_of_quadratic_constraints}
The root-mean-square, total and individual harmonic distortion voltage limits and the root-mean-square current limit are the second-order cone constraints~(47)--(49) and~(51) in~\cite{vanacker_hosting_2026}. Consider a generic constraint $\RVX^{2} + \RVY^{2} \leq a^{2}$ with deterministic~$a$. Following~\cite{vanacker_ivr}, an auxiliary lifted variable~$\RVZ$ is introduced, $\RVX^{2} + \RVY^{2} = \RVZ \leq a^{2}$. Supported by~\eqref{eq:gp_product}, the lifting yields~$|\PCESet|$ constraints,
\begin{IEEEeqnarray}{ l }
    \sum_{\bar{\kappa} \in P^{*}(\PCESet,2)} \mkern-15mu M_{\bar{\kappa},\PCEIndex} \,(\RVX_{\kappa_{1}} \RVX_{\kappa_{2}} + \RVY_{\kappa_{1}} \RVY_{\kappa_{2}}) = \RVZ_{\PCEIndex}, \quad \forall \PCEIndex \in \PCESet. \label{eq:soc_aux_eq}
\end{IEEEeqnarray}
For~$\PCEIndex = 0$ the quadratic form is positive semidefinite, but for~$\PCEIndex \neq 0$ it is indefinite and the equality must be retained: relaxing it would allow~$\RVZ_k = 0$ where the true coefficient is negative, which underestimates the variance and makes the chance constraint over-optimistic. The right-hand side becomes the chance constraint,
\begin{IEEEeqnarray}{ l }
    \RVZ_{0} + \lambda(\varepsilon) \sqrt{ \sum_{\PCEIndex \in \PCESet \backslash \{0\}} \mkern-10mu \langle \Polynoom_{\PCEIndex}, \Polynoom_{\PCEIndex} \rangle \, \RVZ_{\PCEIndex}^{2}} \leq a^{2}, \label{eq:cc_soc}
\end{IEEEeqnarray}
a second-order cone in~$\{\RVZ_{\PCEIndex}\}_{\PCEIndex\in\PCESet}$. With the equalities~\eqref{eq:soc_aux_eq} for~$\PCEIndex \neq 0$, the general stochastic harmonic hosting capacity problem is a non-convex quadratically constrained quadratic program.

\subsection{Reduction to a Second-Order Cone Program}
\label{sec:absolute_equality_socp}

The general problem reduces to a second-order cone program for every fairness criterion that fixes the ratios between the allocations of the harmonic units, i.e., $\CurrentMagnitude_{\Unit,\Harmonic,0} = w_{\Unit}\,\HostingCapacity_{\Harmonic}$ with fixed positive weights~$w_{\Unit}$ and a single hosting capacity~$\HostingCapacity_{\Harmonic}$ per harmonic. Absolute equality, $w_{\Unit} = 1$, and weighted equality, e.g., with weights proportional to the rated power of the units, are such criteria. The Kalai-Smorodinsky solution~(37)--(39) in~\cite{vanacker_hosting_2026} is one when all ratios to the individual maxima are equal at the optimum, with~$w_{\Unit}$ the individual maximum of unit~$\Unit$, and maximin~(35) in~\cite{vanacker_hosting_2026} coincides with absolute equality when raising the current of a unit cannot relax a binding constraint. Criteria with a free allocation, e.g., maximum efficiency or proportional fairness~\cite{chen_guide_2023}, require a current per unit and lead to a semidefinite relaxation instead.

For fixed weights, \eqref{eq:real_harmonic_unit_current_pce}--\eqref{eq:imag_harmonic_unit_current_pce} express every injected current coefficient as a known parameter times~$\HostingCapacity_{\Harmonic}$. As the network equations are linear in the coefficients, all voltage coefficients are linear in~$\HostingCapacity_{\Harmonic}$ as well,
\begin{IEEEeqnarray}{ l }
    \RealVoltage_{\Node,\Harmonic,\PCEIndex} = C^{\text{re}}_{\Node,\Harmonic,\PCEIndex}\, \HostingCapacity_{\Harmonic}, \qquad
    \ImaginaryVoltage_{\Node,\Harmonic,\PCEIndex} = C^{\text{im}}_{\Node,\Harmonic,\PCEIndex}\, \HostingCapacity_{\Harmonic},
    \label{eq:voltage_linear}
\end{IEEEeqnarray}
where~$C^{\text{re}}$ and~$C^{\text{im}}$ follow from the harmonic transfer impedances~$Z_{\Harmonic}$ of the network, obtained by inverting the harmonic admittance matrix, the weights~$w_{\Unit}$ and the trigonometric coefficients~\eqref{eq:trig_quadrature}; branch currents follow analogously. The transformer vector-group enters the admittance matrix with the sign of the sequence of the harmonic. Zero-sequence harmonics additionally require the winding configuration and earthing of the transformers. The lifting then collapses to the precomputed scalars,
\begin{IEEEeqnarray}{ l }
\alpha_{\Node,\Harmonic,\PCEIndex} = \mkern-18mu \sum_{\bar{\kappa} \in P^{*}(\PCESet,2)} \mkern-18mu M_{\bar{\kappa},\PCEIndex} \big( C^{\text{re}}_{\Node,\Harmonic,\bar{\kappa}_{1}} C^{\text{re}}_{\Node,\Harmonic,\bar{\kappa}_{2}} + C^{\text{im}}_{\Node,\Harmonic,\bar{\kappa}_{1}} C^{\text{im}}_{\Node,\Harmonic,\bar{\kappa}_{2}} \big),
\end{IEEEeqnarray}
i.e., the coefficients of~$|\ComplexVoltage_{\Node,\Harmonic}|^2/\HostingCapacity_{\Harmonic}^2$. With~$\LiftedHostingCapacity_{\Harmonic} \geq \HostingCapacity_{\Harmonic}^{2}$, a rotated second-order cone, every chance constraint on a sum of squared harmonic quantities over a set~$\mathcal{H}_{c} \subseteq \HarmonicSet$ of harmonics becomes,
\begin{IEEEeqnarray}{ l }
    \sum_{\Harmonic \in \mathcal{H}_{c}} \alpha_{\Node,\Harmonic,0}\, \LiftedHostingCapacity_{\Harmonic} \leq a_{\Node,c} \nonumber \\
    \quad -\: \lambda(\varepsilon) \sqrt{ \sum_{\PCEIndex \in \PCESet \backslash \{0\}} \mkern-10mu \langle \Polynoom_{\PCEIndex}, \Polynoom_{\PCEIndex} \rangle \Big( \sum_{\Harmonic \in \mathcal{H}_{c}} \alpha_{\Node,\Harmonic,\PCEIndex}\, \LiftedHostingCapacity_{\Harmonic} \Big)^{2}}, \label{eq:socp_cc}
\end{IEEEeqnarray}
a second-order cone in~$\bar{\LiftedHostingCapacity}$: $\mathcal{H}_{c} = \{\Harmonic\}$ for the individual harmonic voltage distortion, which reduces to a linear bound on~$\LiftedHostingCapacity_{\Harmonic}$, and $\mathcal{H}_{c} = \HarmonicSet$ for the total harmonic distortion and the root-mean-square voltage and current limits, where~$a_{\Node,c}$ is the limit net of the fundamental. The stochastic harmonic hosting capacity is thus the second-order cone program,
\begin{IEEEeqnarray}{ l l l }
    \max_{\bar{\HostingCapacity},\,\bar{\LiftedHostingCapacity}} & \sum_{\Harmonic \in \HarmonicSet} \Big( \sum_{\Unit \in \HarmonicUnitSet} w_{\Unit} \Big) \HostingCapacity_{\Harmonic} & \nonumber \\
    \text{subject to} \quad & \HostingCapacity_{\Harmonic}^{2} \leq \LiftedHostingCapacity_{\Harmonic}, & \forall \Harmonic \in \HarmonicSet, \nonumber \\
    & w_{\Unit}\,\HostingCapacity_{\Harmonic} \leq \RatedCurrent_{\Unit}, & \forall \Unit \in \HarmonicUnitSet, \Harmonic \in \HarmonicSet, \nonumber \\
    & w_{\Unit}^{2}\,\LiftedHostingCapacity_{\Harmonic} \leq (\RatedCurrent_{\Unit})^{2}, \quad & \forall \Unit \in \HarmonicUnitSet, \Harmonic \in \HarmonicSet, \nonumber \\
    & \eqref{eq:socp_cc}, & \label{eq:socp}
\end{IEEEeqnarray}
with one pair of variables per harmonic, whatever the size of the network or the cardinality of the basis; the latter only enters the precomputed scalars~$\alpha$. The relaxation~$\HostingCapacity_{\Harmonic}^{2} \leq \LiftedHostingCapacity_{\Harmonic}$ is exact: the chance constraints~\eqref{eq:socp_cc} involve~$\bar{\LiftedHostingCapacity}$ only and the objective is strictly increasing in each~$\HostingCapacity_{\Harmonic}$, so any optimum satisfies~$\HostingCapacity_{\Harmonic} = \sqrt{\LiftedHostingCapacity_{\Harmonic}}$. The rating is therefore imposed on~$\LiftedHostingCapacity_{\Harmonic}$ as well. Since the left-hand side of~\eqref{eq:socp_cc} is not monotone in~$\bar{\LiftedHostingCapacity}$, i.e., under recombination an increase of~$\LiftedHostingCapacity_{\Harmonic}$ may reduce the variance more than it increases the expectation, a bound on~$\HostingCapacity_{\Harmonic}$ alone would allow~$\LiftedHostingCapacity_{\Harmonic} > \HostingCapacity_{\Harmonic}^{2}$ once the rating binds, and the reported hosting capacity would violate the chance constraints. A limit that the fundamental already exhausts, $a_{\Node,c} \leq 0$, leaves no harmonic headroom and forces the contributing~$\LiftedHostingCapacity_{\Harmonic}$ to zero. When every angle range collapses to a point, all variances vanish and~\eqref{eq:socp} reduces exactly to the deterministic model.

